% intro

\section{Motivation}\label{sec:motivation}

Climate mitigation is not always linear. Even when pathways share the same
2050 net-zero target, emissions reductions can be front-loaded, linear, or
delayed, reflecting different policy choices and technology portfolios.

Air pollution and its long-term health impacts are closely tied to the pace of
decarbonization. As a result, cumulative health impacts can be path-dependent,
yet limited research has systematically studied how the nonlinear curvature of
mitigation pathways shapes health outcomes.

Korea provides a timely test case: its emerging 2035--2050 policy framework
explicitly considers both a linear pathway and an early-reduction pathway.
\note{cite policy framework}

\section{Research Question}\label{sec:rq}

How does the curvature of decarbonization pathways affect cumulative health
outcomes on the way to a net-zero target?

\paragraph{Test case.} \PM-attributable deaths under Korea's 2050 carbon
neutrality target.

\paragraph{Hypothesis.} Compared with a linear pathway, an early-reduction
pathway will avoid more cumulative \PM-related deaths by bringing forward
reductions in fossil-fuel-related air-pollutant emissions through earlier
climate mitigation.

\section{Contributions}\label{sec:contributions}

\paragraph{Science.} Trade-offs between scenarios. \discuss{overshoot?}

\paragraph{Analytical.} One of the first studies connecting energy, air
quality, and health modeling. \note{Integrated energy--air quality--health
studies already exist; consider narrowing the claim, e.g.\ to pathway
curvature or to Korea.}

\paragraph{Policy.} The pathways compared are those written into Korean law.

\subsection{Proposed Contributions}\label{sec:proposed}
\note{Proposed by Yehyun}

\paragraph{Modeling.}
\begin{itemize}
  \item A GCAM-KAIST air pollution cluster, mirroring GCAM-USA.
  \item Global InMAP calibrated for Korea.
  \item A facility retirement ordering algorithm (e.g.\ for power plant units)
        based on age, emissions, most recent historical activity, and other
        characteristics. See Section~\ref{sec:retirement}.
\end{itemize}

\paragraph{Data.}
\begin{itemize}
  \item A harmonized dataset of emission factors by sector in Korea.
  \item A cleaned and harmonized historical air quality dataset for Korea.
  \item A harmonized spatial allocation of air pollutant emissions by sector,
        derived from CAPSS and other government sources.
  \item A harmonized dataset of Korean demographics: district-level
        population, age structure, baseline mortality rates, and established
        concentration--response functions for BenMAP-style calculations.
  \item An original, harmonized power plant dataset: monthly generation and
        raw mass emissions, combined with characteristics including location,
        operation dates, fuel type, technology, and stack height. This is
        important for modeling power sector activity.
\end{itemize}

\paragraph{Results.}
\begin{itemize}
  \item A computationally inexpensive end-to-end model that can generate
        scenario outputs with limited GPUs, suited to small university
        research teams. The advantage comes mainly from InMAP as the
        atmospheric model.
\end{itemize}
